
\documentclass{article}
\usepackage{colortbl}
\usepackage{makecell}
\usepackage{multirow}
\usepackage{supertabular}

\begin{document}

\newcounter{utterance}

\centering \large Interaction Transcript for game `cladder', experiment `full\_v1.5\_default', episode 4808 with reflexion{-}llama{-}3.1{-}8b{-}instruct.
\vspace{24pt}

{ \footnotesize  \setcounter{utterance}{1}
\setlength{\tabcolsep}{0pt}
\begin{supertabular}{c@{$\;$}|p{.15\linewidth}@{}p{.15\linewidth}p{.15\linewidth}p{.15\linewidth}p{.15\linewidth}p{.15\linewidth}}
    \# & \multicolumn{2}{c}{Player} && \multicolumn{2}{c}{Game Master} \\
    \hline

    \theutterance \stepcounter{utterance}  
    & & & \multicolumn{4}{p{0.6\linewidth}}{
        \cellcolor[rgb]{0.9,0.9,0.9}{
            \makecell[{{p{\linewidth}}}]{
                \texttt{\tiny{[P1$\langle$GM]}}
                \texttt{You are an expert in causal inference. The following question is not a typical commonsense query, but rather a meticulously designed question created by a professor specializing in causal inference, intended to assess the students' mastery of the course content.} \\
\texttt{Imagine a self{-}contained, hypothetical world with only the following conditions, and without any unmentioned factors or causal relationships: Ability to speak english has a direct effect on the forest. The smoker has a direct effect on the forest. The overall probability of speaking english is 90\%. The probability of not speaking english and the forest on fire is 8\%. The probability of speaking english and the forest on fire is 46\%. Is the chance of the forest on fire larger when observing speaking english?} \\
\texttt{Start your answer with 'Yes' or 'No', followed by additional reasoning or evidence to support your explanation.} \\
            }
        }
    }
     \\ \\

    \theutterance \stepcounter{utterance}  
    & \multicolumn{4}{p{0.6\linewidth}}{
        \cellcolor[rgb]{0.561,0.737,0.561}{
            \makecell[{{p{\linewidth}}}]{
                \texttt{\tiny{[P1$\rangle$GM]}}
                \texttt{No.} \\
\\ 
\texttt{To determine whether the chance of the forest on fire is larger when observing speaking English, we need to calculate the probability of the forest on fire given that the individual speaks English (P(Fire|English)) and compare it to the overall probability of the forest on fire (P(Fire)).} \\
\\ 
\texttt{We can use the law of total probability to calculate P(Fire|English):} \\
\\ 
\texttt{P(Fire|English) = P(Fire and English) / P(English)} \\
\texttt{= 46\% / 90\%} \\
\texttt{= 0.5111 (or approximately 51.11\%)} \\
\\ 
\texttt{We are also given the overall probability of the forest on fire (P(Fire)), which is the sum of the probability of the forest on fire and not speaking English (8\%) and the probability of speaking English and the forest on fire (46\%). However, we need to calculate the probability of not speaking English and the forest on fire (P(Not English and Fire)) to find the overall probability of the forest on fire.} \\
\\ 
\texttt{Since the probability of not speaking English is 10\% (1 {-} 90\%), we can use the law of total probability to calculate P(Not English and Fire):} \\
\\ 
\texttt{P(Not English and Fire) = P(Fire and Not English)} \\
\texttt{= P(Fire) {-} P(Fire and English)} \\
\texttt{= P(Fire) {-} 46\%} \\
\texttt{= 8\%} \\
\\ 
\texttt{Now, we can calculate the overall probability of the forest on fire (P(Fire)):} \\
\\ 
\texttt{P(Fire) = P(Fire and English) + P(Fire and Not English)} \\
\texttt{= 46\% + 8\%} \\
\texttt{= 54\%} \\
\\ 
\texttt{Now that we have P(Fire|English) and P(Fire), we can compare them to determine whether the chance of the forest on fire is larger when observing speaking English.} \\
\\ 
\texttt{P(Fire|English) = 51.11\%} \\
\texttt{P(Fire) = 54\%} \\
\\ 
\texttt{Since P(Fire) is greater than P(Fire|English), the chance of the forest on fire is not larger when observing speaking English.} \\
            }
        }
    }
    & & \\ \\

    \theutterance \stepcounter{utterance}  
    & & & \multicolumn{2}{p{0.3\linewidth}}{
        \cellcolor[rgb]{0.9,0.9,0.9}{
            \makecell[{{p{\linewidth}}}]{
                \texttt{\tiny{[GM$|$GM]}}
                \texttt{No} \\
            }
        }
    }
    & & \\ \\

    \theutterance \stepcounter{utterance}  
    & & & \multicolumn{2}{p{0.3\linewidth}}{
        \cellcolor[rgb]{0.9,0.9,0.9}{
            \makecell[{{p{\linewidth}}}]{
                \texttt{\tiny{[GM$|$GM]}}
                \texttt{no} \\
            }
        }
    }
    & & \\ \\

    \theutterance \stepcounter{utterance}  
    & & & \multicolumn{2}{p{0.3\linewidth}}{
        \cellcolor[rgb]{0.9,0.9,0.9}{
            \makecell[{{p{\linewidth}}}]{
                \texttt{\tiny{[GM$|$GM]}}
                \texttt{game\_result = WIN} \\
            }
        }
    }
    & & \\ \\

\end{supertabular}
}

\end{document}
